% GENERATED FILE. Edit canonical lessons, then run npm run generate:books.
% canonical-source-hash: fnv1a64:154ddffdece18fd0
% canonical-lessons: KA-C515-kuruda, KA-C515-kivuda, KA-C515-muka, KA-C515-naralu, KA-C515-nalugu

\chapter{Blindness, Deafness and Pain}
\label{ch:ka-blindness-deafness-and-pain}
\hlchaptermodality{515}

\begin{chapteropening}
\textbf{By the end of this chapter:} \emph{I can say a blind man, deaf, mute, to groan and to droop.}

Complete the last lesson of chapter 515: I can say a blind man, deaf, mute, to groan and to droop.
\end{chapteropening}

\section[\texorpdfstring{\kn{ಕುರುಡ} (kuruḍa)}{ಕುರುಡ (kuruḍa)}]{\kn{ಕುರುಡ} (kuruḍa) — a blind man; blind}

\label{lesson:KA-C515-kuruda}



\begin{quote}
Before the new one: say the Kannada for to get tired, then the Kannada for lifespan.
\end{quote}

\subsection*{You'll want to know: \kn{ಕುರುಡ}}
\textbf{\kn{ಕುರುಡ}} — \emph{kuruḍa} — \textquotedblleft{}a blind man; blind\textquotedblright{}.

\begin{quote}
\kn{ಆ} \kn{ಕುರುಡ} \kn{ಸಂಗೀತಗಾರ} \kn{ಸುಂದರವಾಗಿ} \kn{ಹಾಡುತ್ತಾನೆ}.
\end{quote}

\begin{quote}
\emph{ā kuruḍa saṅgītagāra sundaravāgi hāḍuttāne.}
\end{quote}

\begin{quote}
That blind musician sings beautifully.
\end{quote}

Plural: \kn{ಕುರುಡರು}, \emph{kuruḍaru}.

A blind woman is \kn{ಕುರುಡಿ}. Politer: \kn{ದೃಷ್ಟಿ} \kn{ಇಲ್ಲದವರು}, those without sight.

\begin{cousinweb}
\textbf{\kn{ಕುರುಡ}} (\emph{kuruḍa}) is native Dravidian: DEDR 1787 lists Kannada \emph{kuruḍu}, \textquotedblleft{}blindness\textquotedblright{}, and \emph{kuruḍa}, \textquotedblleft{}a blind man\textquotedblright{}.
\end{cousinweb}

\begin{grammarlens}[title={What you've built}]
One more word for describing what you lived and felt.
\end{grammarlens}

\subsection*{Your turn}
\begin{itemize}
\raggedright
  \item \emph{Say it:} \emph{kuruḍa}
  \item \emph{Say it:} \emph{kuruḍa}, once more
  \item \emph{Say it:} \emph{ā kuruḍa saṅgītagāra sundaravāgi hāḍuttāne.} — That blind musician sings beautifully.
  \item \emph{Recall:} say the Kannada for to get tired, then the Kannada for lifespan, then say \emph{kuruḍa} again
\end{itemize}

\begin{tcolorbox}[breakable,colback=teal!4,colframe=teal!35!black,title={Before you move on}]
What does \kn{ಕುರುಡ} mean? (\textquotedblleft{}A blind man; blind\textquotedblright{}.) Say it once more.
\end{tcolorbox}
\section[\texorpdfstring{\kn{ಕಿವುಡ} (kivuḍa)}{ಕಿವುಡ (kivuḍa)}]{\kn{ಕಿವುಡ} (kivuḍa) — deaf; a deaf man}

\label{lesson:KA-C515-kivuda}



\begin{quote}
Before the new one: say the Kannada for lifespan, then the Kannada for a blind man.
\end{quote}

\subsection*{You'll want to know: \kn{ಕಿವುಡ}}
\textbf{\kn{ಕಿವುಡ}} — \emph{kivuḍa} — \textquotedblleft{}deaf; a deaf man\textquotedblright{}.

\begin{quote}
\kn{ಅಜ್ಜ} \kn{ಸ್ವಲ್ಪ} \kn{ಕಿವುಡ}, \kn{ಜೋರಾಗಿ} \kn{ಮಾತಾಡಬೇಕು}.
\end{quote}

\begin{quote}
\emph{ajja svalpa kivuḍa, jōrāgi mātāḍabēku.}
\end{quote}

\begin{quote}
Grandfather is a little deaf; you have to speak loudly.
\end{quote}

Built on \kn{ಕಿವಿ}, the ear. A deaf woman is \kn{ಕಿವುಡಿ}.

\begin{cousinweb}
\textbf{\kn{ಕಿವುಡ}} (\emph{kivuḍa}) is native Dravidian: DEDR 1977 lists Kannada \emph{kivi}, \textquotedblleft{}ear\textquotedblright{}, and \emph{kivuḍa}, \textquotedblleft{}a deaf man\textquotedblright{}.
\end{cousinweb}

\begin{grammarlens}[title={What you've built}]
One more word for describing what you lived and felt.
\end{grammarlens}

\subsection*{Your turn}
\begin{itemize}
\raggedright
  \item \emph{Say it:} \emph{kivuḍa}
  \item \emph{Say it:} \emph{kivuḍa}, once more
  \item \emph{Say it:} \emph{ajja svalpa kivuḍa, jōrāgi mātāḍabēku.} — Grandfather is a little deaf; you have to speak loudly.
  \item \emph{Recall:} say the Kannada for lifespan, then the Kannada for a blind man, then say \emph{kivuḍa} again
\end{itemize}

\begin{tcolorbox}[breakable,colback=teal!4,colframe=teal!35!black,title={Before you move on}]
What does \kn{ಕಿವುಡ} mean? (\textquotedblleft{}Deaf; a deaf man\textquotedblright{}.) Say it once more.
\end{tcolorbox}
\section[\texorpdfstring{\kn{ಮೂಕ} (mūka)}{ಮೂಕ (mūka)}]{\kn{ಮೂಕ} (mūka) — mute; unable to speak}

\label{lesson:KA-C515-muka}



\begin{quote}
Before the new one: say the Kannada for a blind man, then the Kannada for deaf.
\end{quote}

\subsection*{You'll want to know: \kn{ಮೂಕ}}
\textbf{\kn{ಮೂಕ}} — \emph{mūka} — \textquotedblleft{}mute; unable to speak\textquotedblright{}.

\begin{quote}
\kn{ಆ} \kn{ಮೂಕ} \kn{ಹುಡುಗ} \kn{ಕೈಸನ್ನೆಯಿಂದ} \kn{ಮಾತಾಡುತ್ತಾನೆ}.
\end{quote}

\begin{quote}
\emph{ā mūka huḍuga kaisanneyinda mātāḍuttāne.}
\end{quote}

\begin{quote}
That mute boy talks with hand signs.
\end{quote}

\kn{ಮೂಕ} \kn{ಪ್ರಾಣಿ}, a dumb creature, an animal that cannot complain.

\begin{cousinweb}
\textbf{\kn{ಮೂಕ}} (\emph{mūka}) is borrowed from Sanskrit \emph{mūka}, \textquotedblleft{}tied, tongue-tied, dumb, silent\textquotedblright{} (Monier-Williams).
\end{cousinweb}

\begin{grammarlens}[title={What you've built}]
One more word for describing what you lived and felt.
\end{grammarlens}

\subsection*{Your turn}
\begin{itemize}
\raggedright
  \item \emph{Say it:} \emph{mūka}
  \item \emph{Say it:} \emph{mūka}, once more
  \item \emph{Say it:} \emph{ā mūka huḍuga kaisanneyinda mātāḍuttāne.} — That mute boy talks with hand signs.
  \item \emph{Recall:} say the Kannada for a blind man, then the Kannada for deaf, then say \emph{mūka} again
\end{itemize}

\begin{tcolorbox}[breakable,colback=teal!4,colframe=teal!35!black,title={Before you move on}]
What does \kn{ಮೂಕ} mean? (\textquotedblleft{}Mute; unable to speak\textquotedblright{}.) Say it once more.
\end{tcolorbox}
\section[\texorpdfstring{\kn{ನರಳು} (naraḷu)}{ನರಳು (naraḷu)}]{\kn{ನರಳು} (naraḷu) — to groan; to moan}

\label{lesson:KA-C515-naralu}



\begin{quote}
Before the new one: say the Kannada for deaf, then the Kannada for mute.
\end{quote}

\subsection*{You'll want to know: \kn{ನರಳು}}
\textbf{\kn{ನರಳು}} — \emph{naraḷu} — \textquotedblleft{}to groan; to moan\textquotedblright{}.

\begin{quote}
\kn{ನೋವಿನಿಂದ} \kn{ರೋಗಿ} \kn{ರಾತ್ರಿಯೆಲ್ಲ} \kn{ನರಳಿದನು}.
\end{quote}

\begin{quote}
\emph{nōvininda rōgi rātriyella naraḷidanu.}
\end{quote}

\begin{quote}
The patient groaned with pain all night.
\end{quote}

Past (I …): \kn{ನರಳಿದೆ}, \emph{naraḷide}.

A low sound of pain.

\begin{cousinweb}
\textbf{\kn{ನರಳು}} (\emph{naraḷu}) is native Dravidian: DEDR 2904 lists Kannada \emph{naral} and \emph{neraḷu}, \textquotedblleft{}to groan, moan\textquotedblright{}.
\end{cousinweb}

\begin{grammarlens}[title={What you've built}]
One more word for describing what you lived and felt.
\end{grammarlens}

\subsection*{Your turn}
\begin{itemize}
\raggedright
  \item \emph{Say it:} \emph{naraḷu}
  \item \emph{Say it:} \emph{naraḷu}, once more
  \item \emph{Say it:} \emph{nōvininda rōgi rātriyella naraḷidanu.} — The patient groaned with pain all night.
  \item \emph{Recall:} say the Kannada for deaf, then the Kannada for mute, then say \emph{naraḷu} again
\end{itemize}

\begin{tcolorbox}[breakable,colback=teal!4,colframe=teal!35!black,title={Before you move on}]
What does \kn{ನರಳು} mean? (\textquotedblleft{}To groan; to moan\textquotedblright{}.) Say it once more.
\end{tcolorbox}
\section[\texorpdfstring{\kn{ನಲುಗು} (nalugu)}{ನಲುಗು (nalugu)}]{\kn{ನಲುಗು} (nalugu) — to droop; to be worn down}

\label{lesson:KA-C515-nalugu}



\begin{quote}
Before the new one: say the Kannada for mute, then the Kannada for to groan.
\end{quote}

\subsection*{You'll want to know: \kn{ನಲುಗು}}
\textbf{\kn{ನಲುಗು}} — \emph{nalugu} — \textquotedblleft{}to droop; to be worn down\textquotedblright{}.

\begin{quote}
\kn{ಬಿಸಿಲಿಗೆ} \kn{ತೋಟದ} \kn{ಹೂವುಗಳು} \kn{ನಲುಗಿದವು}.
\end{quote}

\begin{quote}
\emph{bisilige tōṭada hūvugaḷu nalugidavu.}
\end{quote}

\begin{quote}
The garden flowers drooped in the sun.
\end{quote}

Past (I …): \kn{ನಲುಗಿದೆ}, \emph{nalugide}.

Of people too: \kn{ಕಷ್ಟದಿಂದ} \kn{ನಲುಗು}, to be worn down by hardship.

\begin{cousinweb}
\textbf{\kn{ನಲುಗು}} (\emph{nalugu}) is native Dravidian: DEDR 3611 lists Kannada \emph{nalagu}, \emph{nalugu}, \textquotedblleft{}to become rumpled, to droop, to grow weak or faint from sickness, trouble or grief\textquotedblright{}.
\end{cousinweb}

\begin{grammarlens}[title={What you've built}]
One more word for describing what you lived and felt.
\end{grammarlens}

\subsection*{Your turn}
\begin{itemize}
\raggedright
  \item \emph{Say it:} \emph{nalugu}
  \item \emph{Say it:} \emph{nalugu}, once more
  \item \emph{Say it:} \emph{bisilige tōṭada hūvugaḷu nalugidavu.} — The garden flowers drooped in the sun.
  \item \emph{Recall:} say the Kannada for mute, then the Kannada for to groan, then say \emph{nalugu} again
\end{itemize}

\begin{tcolorbox}[breakable,colback=teal!4,colframe=teal!35!black,title={Before you move on}]
What does \kn{ನಲುಗು} mean? (\textquotedblleft{}To droop; to be worn down\textquotedblright{}.) Say it once more.
\end{tcolorbox}
